\documentclass[10pt,a4paper]{amsart}
\usepackage{amsmath,amssymb,amsthm,graphicx}
\usepackage{verbatim}
\newenvironment{reference}{\comment}{\endcomment}
%\newenvironment{reference}{}{}
\newenvironment{slogan}{\comment}{\endcomment}
\newenvironment{history}{\comment}{\endcomment}

% For commutative diagrams we use Xy-pic
\usepackage[all]{xy}

% We use 2cell for 2-commutative diagrams.
\xyoption{2cell}
% Theorem environments.
%
\theoremstyle{plain}
\newtheorem{theorem}[subsection]{Theorem}
\newtheorem{proposition}[subsection]{Proposition}
\newtheorem{lemma}[subsection]{Lemma}

\theoremstyle{definition}
\newtheorem{definition}[subsection]{Definition}
\newtheorem{example}[subsection]{Example}
\newtheorem{exercise}[subsection]{Exercise}
\newtheorem{situation}[subsection]{Situation}

\theoremstyle{remark}
\newtheorem{remark}[subsection]{Remark}
\newtheorem{remarks}[subsection]{Remarks}

\numberwithin{equation}{subsection}

% Macros
%
\def\lim{\mathop{\mathrm{lim}}\nolimits}
\def\colim{\mathop{\mathrm{colim}}\nolimits}
\def\Spec{\mathop{\mathrm{Spec}}}
\def\Hom{\mathop{\mathrm{Hom}}\nolimits}
\def\Ext{\mathop{\mathrm{Ext}}\nolimits}
\def\SheafHom{\mathop{\mathcal{H}\!\mathit{om}}\nolimits}
\def\SheafExt{\mathop{\mathcal{E}\!\mathit{xt}}\nolimits}
\def\Sch{\mathit{Sch}}
\def\Mor{\mathop{\mathrm{Mor}}\nolimits}
\def\Ob{\mathop{\mathrm{Ob}}\nolimits}
\def\Sh{\mathop{\mathit{Sh}}\nolimits}
\def\NL{\mathop{N\!L}\nolimits}
\def\CH{\mathop{\mathrm{CH}}\nolimits}
\def\proetale{{pro\text{-}\acute{e}tale}}
\def\etale{{\acute{e}tale}}
\def\QCoh{\mathit{QCoh}}
\def\Ker{\mathop{\mathrm{Ker}}}
\def\Im{\mathop{\mathrm{Im}}}
\def\Coker{\mathop{\mathrm{Coker}}}
\def\Coim{\mathop{\mathrm{Coim}}}

% Boxtimes
%
\DeclareMathSymbol{\boxtimes}{\mathbin}{AMSa}{"02}

%
% Macros for moduli stacks/spaces
%
\def\QCohstack{\mathcal{QC}\!\mathit{oh}}
\def\Cohstack{\mathcal{C}\!\mathit{oh}}
\def\Spacesstack{\mathcal{S}\!\mathit{paces}}
\def\Quotfunctor{\mathrm{Quot}}
\def\Hilbfunctor{\mathrm{Hilb}}
\def\Curvesstack{\mathcal{C}\!\mathit{urves}}
\def\Polarizedstack{\mathcal{P}\!\mathit{olarized}}
\def\Complexesstack{\mathcal{C}\!\mathit{omplexes}}
% \Pic is the operator that assigns to X its picard group, usage \Pic(X)
% \Picardstack_{X/B} denotes the Picard stack of X over B
% \Picardfunctor_{X/B} denotes the Picard functor of X over B
\def\Pic{\mathop{\mathrm{Pic}}\nolimits}
\def\Picardstack{\mathcal{P}\!\mathit{ic}}
\def\Picardfunctor{\mathrm{Pic}}
\def\Deformationcategory{\mathcal{D}\!\mathit{ef}}
\usepackage{cleveref}
\crefname{originallemma}{original lemma}{original lemmas}
\Crefname{originallemma}{Original lemma}{Original lemmas}
\makeatletter
\newlabel{lemma-lisse-etale}{{0788}{source}{}{sourcepointer-1}{}}
\newlabel{lemma-lisse-etale@cref}{{[originallemma][0][]0788}{[1][0][]source}{}{}{}}
\makeatother
\crefname{originallemma}{original lemma}{original lemmas}
\Crefname{originallemma}{Original lemma}{Original lemmas}
\makeatletter
\newlabel{lemma-lisse-etale-modules}{{0789}{source}{}{sourcepointer-2}{}}
\newlabel{lemma-lisse-etale-modules@cref}{{[originallemma][0][]0789}{[1][0][]source}{}{}{}}
\makeatother
\crefname{originalreference}{original reference}{original references}
\Crefname{originalreference}{Original reference}{Original references}
\makeatletter
\newlabel{sites-modules-definition-g-shriek}{{04BF}{source}{}{sourcepointer-3}{}}
\newlabel{sites-modules-definition-g-shriek@cref}{{[originalreference][0][]04BF}{[1][0][]source}{}{}{}}
\makeatother
\crefname{originalreference}{original reference}{original references}
\Crefname{originalreference}{Original reference}{Original references}
\makeatletter
\newlabel{spaces-topologies-lemma-etale-dominates-smooth}{{0CFV}{source}{}{sourcepointer-4}{}}
\newlabel{spaces-topologies-lemma-etale-dominates-smooth@cref}{{[originalreference][0][]0CFV}{[1][0][]source}{}{}{}}
\makeatother
\title{The structure sheaf of the lisse-etale and flat-fppf sites pushes forward to the structure sheaf}
\author{The Stacks Project Authors}
\date{stacks-project:stacks-cohomology.tex:local-snapshot-20261012; GNU FDL 1.2 or later}
\begin{document}
\maketitle
\noindent\textit{Editorial note.} All delivered mathematical text is the author's own, copied verbatim from the pinned Stacks Project chapter file stacks-cohomology.tex; no statement or proof step was rewritten, reordered or removed.\par\medskip
\noindent\textit{Editorial note.} The counted span cites the lemmas recorded in assembly.external\_references with canonical tags as navigation pointers.\par\medskip
\noindent\textit{Allowed standard prerequisites (editorial).}\begin{itemize}
\item The g\_! functor. Modules on Sites, Definition sites-modules-definition-g-shriek. Colimit presheaf formula for g\_!O'.
\item Etale dominates smooth. Topologies on Spaces, Lemma spaces-topologies-lemma-etale-dominates-smooth. Lets V be replaced so that V -> U exists.
\end{itemize}\medskip
\ifdefined\hypertarget\hypertarget{sourcepointer-1}{}\fi
\noindent\textit{Source reference.} Original lemma 0788: sources/stacks-chow/tags.txt, source lines 9228--9228. Cross-chapter standard prerequisite cited by the author inside the counted span; it is used as a pointer and its own proof is not part of this item.\par\medskip
\ifdefined\hypertarget\hypertarget{sourcepointer-2}{}\fi
\noindent\textit{Source reference.} Original lemma 0789: sources/stacks-chow/tags.txt, source lines 9229--9229. Cross-chapter standard prerequisite cited by the author inside the counted span; it is used as a pointer and its own proof is not part of this item.\par\medskip
\ifdefined\hypertarget\hypertarget{sourcepointer-3}{}\fi
\noindent\textit{Source reference.} Original reference 04BF: sources/stacks-chow/tags.txt, source lines 5650--5650. Cross-chapter standard prerequisite cited by the author inside the counted span; it is used as a pointer and its own proof is not part of this item.\par\medskip
\ifdefined\hypertarget\hypertarget{sourcepointer-4}{}\fi
\noindent\textit{Source reference.} Original reference 0CFV: sources/stacks-chow/tags.txt, source lines 15625--15625. Cross-chapter standard prerequisite cited by the author inside the counted span; it is used as a pointer and its own proof is not part of this item.\par\medskip
\begin{lemma}
\label{lemma-lisse-etale-structure-sheaf}
Let $\mathcal{X}$ be an algebraic stack. Notation as in
Lemmas \ref{lemma-lisse-etale} and \ref{lemma-lisse-etale-modules}.
\begin{enumerate}
\item We have $g_!\mathcal{O}_{\mathcal{X}_{lisse,\etale}} =
\mathcal{O}_\mathcal{X}$.
\item We have $g_!\mathcal{O}_{\mathcal{X}_{flat, fppf}} =
\mathcal{O}_\mathcal{X}$.
\end{enumerate}
\end{lemma}

\begin{proof}
In this proof we write
$\mathcal{C} = \mathcal{X}_\etale$
(resp.\ $\mathcal{C} = \mathcal{X}_{fppf}$)
and we denote
$\mathcal{C}' = \mathcal{X}_{lisse,\etale}$
(resp.\ $\mathcal{C}' = \mathcal{X}_{flat, fppf}$).
Then $\mathcal{C}'$ is a full subcategory of $\mathcal{C}$.
In this proof we will think of objects $V$ of $\mathcal{C}$
as schemes over $\mathcal{X}$ and objects $U$ of $\mathcal{C}'$
as schemes smooth (resp.\ flat) over $\mathcal{X}$.
Finally, we write $\mathcal{O} = \mathcal{O}_\mathcal{X}$
and $\mathcal{O}' = \mathcal{O}_{\mathcal{X}_{lisse,\etale}}$
(resp.\ $\mathcal{O}' = \mathcal{O}_{\mathcal{X}_{flat,fppf}}$).
In the notation above we have $\mathcal{O}(V) = \Gamma(V, \mathcal{O}_V)$
and $\mathcal{O}'(U) = \Gamma(U, \mathcal{O}_U)$.
Consider the $\mathcal{O}$-module homomorphism
$g_!\mathcal{O}' \to \mathcal{O}$
adjoint to the identification $\mathcal{O}' = g^{-1}\mathcal{O}$.

\medskip\noindent
Recall that $g_!\mathcal{O}'$ is the sheaf associated to the presheaf
$g_{p!}\mathcal{O}'$ given by the rule
$$
V \longmapsto \colim_{V \to U} \mathcal{O}'(U)
$$
where the colimit is taken in the category of abelian groups
(Modules on Sites, Definition \ref{sites-modules-definition-g-shriek}).
Below we will use frequently that if
$$
V \to U \to U'
$$
are morphisms and if $f' \in \mathcal{O}'(U')$ restricts to
$f \in \mathcal{O}'(U)$, then $(V \to U, f)$ and $(V \to U', f')$
define the same element of the colimit. Also,
$g_!\mathcal{O}' \to \mathcal{O}$ maps the element
$(V \to U, f)$ simply to the pullback of $f$ to $V$.

\medskip\noindent
Let us prove that $g_!\mathcal{O}' \to \mathcal{O}$ is surjective.
Let $h \in \mathcal{O}(V)$ for some object $V$ of $\mathcal{C}$.
It suffices to show that $h$ is locally in the image. Choose an object
$U$ of $\mathcal{C}'$ corresponding to a surjective smooth morphism
$U \to \mathcal{X}$. Since $U \times_\mathcal{X} V \to V$ is surjective
smooth, after replacing $V$ by the members of an \'etale covering of $V$
we may assume there exists a morphism $V \to U$, see
Topologies on Spaces, Lemma
\ref{spaces-topologies-lemma-etale-dominates-smooth}. Using $h$ we obtain
a morphism $V \to U \times \mathbf{A}^1$ such that writing
$\mathbf{A}^1 = \Spec(\mathbf{Z}[t])$ the element
$t \in \mathcal{O}(U \times \mathbf{A}^1)$ pulls back to $h$.
Since $U \times \mathbf{A}^1$ is an object of $\mathcal{C}'$
we see that $(V \to U \times \mathbf{A}^1, t)$
is an element of the colimit above which maps to $h \in \mathcal{O}(V)$
as desired.

\medskip\noindent
Suppose that $s \in g_!\mathcal{O}'(V)$ is a section
mapping to zero in $\mathcal{O}(V)$. To finish the proof we have to show
that $s$ is zero. After replacing $V$ by the members
of a covering we may assume $s$ is an element of the colimit
$$
\colim_{V \to U} \mathcal{O}'(U)
$$
Say $s = \sum (\varphi_i, s_i)$ is a finite sum with
$\varphi_i : V \to U_i$, $U_i$ smooth (resp.\ flat) over $\mathcal{X}$, and
$s_i \in \Gamma(U_i, \mathcal{O}_{U_i})$. Choose a scheme $W$ surjective
\'etale over the algebraic space
$U = U_1 \times_\mathcal{X} \ldots \times_\mathcal{X} U_n$.
Note that $W$ is still smooth (resp.\ flat) over $\mathcal{X}$, i.e.,
defines an object of $\mathcal{C}'$. The fibre product
$$
V' = V \times_{(\varphi_1, \ldots, \varphi_n), U} W
$$
is surjective \'etale over $V$, hence it suffices to show that $s$ maps
to zero in $g_!\mathcal{O}'(V')$. Note that the restriction
$\sum (\varphi_i, s_i)|_{V'}$ corresponds to the sum of the pullbacks
of the functions $s_i$ to $W$. In other words, we have reduced to the case
of $(\varphi, s)$ where $\varphi : V \to U$ is a morphism with $U$ in
$\mathcal{C}'$ and $s \in \mathcal{O}'(U)$ restricts to zero in
$\mathcal{O}(V)$. By the commutative diagram
$$
\xymatrix{
V \ar[rr]_-{(\varphi, 0)} \ar[rrd]_\varphi & & U \times \mathbf{A}^1 \\
& & U \ar[u]_{(\text{id}, 0)}
}
$$
we see that
$((\varphi, 0) : V \to U \times \mathbf{A}^1, \text{pr}_2^*x)$
represents zero in the colimit above. Hence we may
replace $U$ by $U \times \mathbf{A}^1$, $\varphi$ by $(\varphi, 0)$
and $s$ by $\text{pr}_1^*s + \text{pr}_2^*x$. Thus we may assume that
the vanishing locus $Z : s = 0$ in $U$ of $s$ is smooth (resp.\ flat)
over $\mathcal{X}$. Then we see that $(V \to Z, 0)$ and $(\varphi, s)$
have the same value in the colimit, i.e., we see that the element $s$
is zero as desired.
\end{proof}

\end{document}
